\section{Producing the angle certificate from CP transport}
\label{urgent:cp}
\noindent\textbf{Status.} This is an unchecked proof expansion of
\texttt{q:lem-schur} and the geometric input to \texttt{q:thm-wire-gap}.
No angle certificate is an assumption of the final producer below.
All spaces in this section are finite-dimensional complex Hilbert spaces.
Inner products are conjugate-linear in the first argument. Empty subspaces
are allowed; an angle involving a zero reduced subspace is zero.

\subsection{Primitive CP-order facts}
A map $\Phi:M_d\to M_e$ is completely positive when all its matrix
amplifications preserve positive semidefinite matrices. Equivalently in
finite dimensions, $\Phi(X)=\sum_a K_aXK_a^*$ for a finite Kraus family.
For two such maps the tensor map has Kraus family $K_a\otimes L_b$:
\[
 (\Phi\otimes\Theta)(X)=\sum_{a,b}(K_a\otimes L_b)X(K_a\otimes L_b)^*.
\]
This proves positivity on entangled inputs; positivity only on elementary
tensors would not suffice. Composition, addition, and multiplication by a
nonnegative real number preserve CP. If $0\le_{\rm CP}A\le_{\rm CP}B$
and $0\le_{\rm CP}C\le_{\rm CP}D$, then
\[
 B\otimes D-A\otimes C=(B-A)\otimes D+A\otimes(D-C)\ge_{\rm CP}0.
\]
Induction gives the corresponding finite tensor-product order statement.
These are foundational CP lemmas, not instances of the desired angle bound.

Suppose $-x\Psi\le_{\rm CP}D\le_{\rm CP}x\Psi$ and
$-y\Theta\le_{\rm CP}E\le_{\rm CP}y\Theta$. The identities
\begin{align*}
xy\Psi\otimes\Theta-D\otimes E
 &=\tfrac12\big[(x\Psi+D)\otimes(y\Theta-E)
                 +(x\Psi-D)\otimes(y\Theta+E)\big],\\
xy\Psi\otimes\Theta+D\otimes E
 &=\tfrac12\big[(x\Psi+D)\otimes(y\Theta+E)
                 +(x\Psi-D)\otimes(y\Theta-E)\big]
\end{align*}
express both sides as CP maps. Thus
\[
 -xy\Psi\otimes\Theta\le_{\rm CP}D\otimes E
                         \le_{\rm CP}xy\Psi\otimes\Theta.
\tag{CP-MIX}
\]
No positivity of either defect $D,E$ is asserted or needed.

\subsection{Strict Gram positivity and injectivity}
Let $C=A^*A\ne0$, let $\tau_j,R_j>0$, and let
$\Psi_j(X)=\operatorname{tr}(\tau_jX)R_j$. Suppose
$(1-\delta_j)\Psi_j\le_{\rm CP}\Phi_j$, with $0\le\delta_j<1$.
Finite tensor order gives
\[
 Q:=(\bigotimes_j\Phi_j)(C)\ \ge
 \left[\prod_j(1-\delta_j)\right]
 \operatorname{tr}[(\bigotimes_j\tau_j)C]\bigotimes_jR_j>0.
\tag{CP-INJ}
\]
For strictness, write $T=\bigotimes_j\tau_j$. It is positive definite,
so $T\ge tI$ for some $t>0$. Since $C\ge0$ and $C\ne0$,
$\operatorname{tr}C>0$, and
$\operatorname{tr}(TC)\ge t\operatorname{tr}C>0$; the inequality follows
by taking the trace of $C^{1/2}(T-tI)C^{1/2}\ge0$.
Every other factor in (CP-INJ) is strictly positive. If $Q=T_v^*T_v$,
then $T_v\xi=0$ implies $\langle\xi,Q\xi\rangle=0$ and hence $\xi=0$.
Tensor products of these injective maps are injective: each has a left
inverse on its image, and their tensor product is a left inverse.
A tensor product of nonzero edge and boundary vectors is nonzero, since
its squared norm is the product of their positive squared norms. This
also proves that the proposed common ground space is nonzero.

The physical Gram identity must be proved in coordinates with the
conjugation convention fixed. If $W_{p,a,b}$ is a wire amplitude, put
$(K_p)_{b,a}=\overline{W_{p,a,b}}$. Direct finite summation gives
\[
 (T_v^*T_v)_{b',b}
 =\sum_{p,a',a}\overline{W_{p,a',b'}}C_{a',a}W_{p,a,b}
 =\sum_p(K_pCK_p^*)_{b',b}.
\]
Iterating over legs gives $T_v^*T_v=(\bigotimes_j\Phi_j)(C)$.
There is no partial transpose of $C$.

\subsection{Quotient metric: a proof with all subspaces specified}
Let $p,q$ be commuting orthogonal projections and put
\[
 W=\operatorname{ran}(pq),\quad U=\operatorname{ran}(p(I-q)),\quad
 V=\operatorname{ran}((I-p)q),\quad Z=\operatorname{ran}((I-p)(I-q)).
\]
The four products are mutually orthogonal self-adjoint projections summing
to $I$. In particular $\operatorname{ran}p=W\oplus U$,
$\operatorname{ran}q=W\oplus V$, and their intersection is $W$.
For an invertible $L$, $L(M\cap N)=LM\cap LN$: apply $L^{-1}$ to both
members of an equality $Lm=Ln$. Therefore the common image under $G^{1/2}$
is exactly $G^{1/2}W$.

Assume $G=I+A+B+D\ge mI>0$, $[A,q]=[B,p]=0$,
$\|A\|\le x$, $\|B\|\le y$, $\|D\|\le xy$, where
$0\le x,y<1$ and $m=(1-x)(1-y)$. For $E,F\in\{U,V,W\}$ set
$G_{EF}=\pi_EG\iota_F$, with inclusion $\iota_F$ and orthogonal
projection $\pi_E$. Compression gives $G_{WW}\ge mI_W$, so it is
invertible on $W$ and $\|G_{WW}^{-1}\|\le m^{-1}$.
For $u\in U$ define $w(u)=-G_{WW}^{-1}G_{WU}u$ and $j_Uu=u+w(u)$.
Then
\[
 \pi_WGj_Uu=0,\qquad
 \langle u+w,G(u+w)\rangle
 =\langle j_Uu,Gj_Uu\rangle+
   \langle w-w(u),G_{WW}(w-w(u))\rangle.
\tag{SCHUR-MIN}
\]
This follows by expanding and using the displayed orthogonality; it proves
both existence and uniqueness of the minimizer. Define $j_V$ identically.
The maps $G^{1/2}j_U$ and $G^{1/2}j_V$ parametrize precisely the two
subspaces after removing $G^{1/2}W$ orthogonally. Indeed, subtracting the
minimizer is necessary and sufficient for orthogonality to that image.
The $Z$ sector is never minimized over: it is not in the intersection.

For $E,F\in\{U,V\}$ the resulting Gram blocks are
\[
 S_{EF}=G_{EF}-G_{EW}G_{WW}^{-1}G_{WF}.
\tag{SCHUR-BLOCK}
\]
For example, expansion of $\langle j_Uu,Gj_Vv\rangle$ gives
$\langle u,S_{UV}v\rangle$. Since $u\perp w(u)$,
\[
 \langle u,S_{UU}u\rangle
 =\langle j_Uu,Gj_Uu\rangle\ge m\|j_Uu\|^2\ge m\|u\|^2,
\]
and likewise $S_{VV}\ge mI$. Thus the maps
$G^{1/2}j_US_{UU}^{-1/2}$ and $G^{1/2}j_VS_{VV}^{-1/2}$ are isometries
onto the reduced subspaces. Their cross-Gram matrix proves, by the
definition of operator norm, that the Friedrichs cosine is exactly
\[
 \|S_{UU}^{-1/2}S_{UV}S_{VV}^{-1/2}\|.
\tag{SCHUR-ANGLE}
\]
This identity is derived, not postulated as an angle certificate.

Commutation implies that $A$ preserves the $q$-eigenspaces and $B$ the
$p$-eigenspaces. The orthogonal blocks therefore satisfy
\[
 G_{UV}=D_{UV},\quad G_{UW}=B_{UW}+D_{UW},\quad
 G_{WV}=A_{WV}+D_{WV}.
\]
Compression is contractive, and hence
\[
 \|S_{UV}\|\le xy+(y+xy)m^{-1}(x+xy).
\]
The spectral theorem applied to $S_{UU},S_{VV}$ gives inverse-square-root
norms at most $m^{-1/2}$. Substitution into (SCHUR-ANGLE) yields
\[
 c_F\le\frac{xy}{m}+\frac{xy(1+x)(1+y)}{m^2}
 =\frac{2xy(1+xy)}{(1-x)^2(1-y)^2}=:\mathcal B(x,y).
\tag{SCHUR-FINAL}
\]
The last equality is ordinary field arithmetic. It does not require any
spectral estimate beyond the stated compression coercivity.

\subsection{Whitening and the actual physical edge subspaces}
Consider distinct edges $e=uv$ and $f=vw$ of a simple graph, so $u\ne w$.
Remove the isometric polar factors of the injective block maps. On the
virtual block spaces the physical maps are the positive invertible filters
$S_a=(T_a^*T_a)^{1/2}$. Let $E_e$ be the virtual subspace which fixes the
edge vector on its two ports and leaves all other ports free; define $E_f$
similarly. The physical two-block regional subspace, within the block
images and with the third block free, is exactly
\[
 L E_e,\quad L E_f,\qquad L=S_u\otimes S_v\otimes S_w.
\tag{PHYS-RANGE}
\]
This equality follows directly from contraction of the edge vector and
surjectivity of every invertible outside filter on a free factor.
It is not an assumption about orthogonal projectors.

Keep all unselected wire maps at their actual values and denote their
tensor product by $\mathcal R$. On the two selected ports set
$D_i=\Phi_i-\Psi_i$, $D_j=\Phi_j-\Psi_j$, and
\[
 G_0=(\Psi_i\otimes\Psi_j\otimes\mathcal R)(C_v)
     =R_i\otimes R_j\otimes R_0>0.
\]
The equality follows by taking the two weighted partial traces in the
rank-one maps. Strictness follows from (CP-INJ); consequently $R_0>0$.
The reference filter
$L_0=S_u\otimes G_0^{1/2}\otimes S_w$ factors on the two selected ports
and the remaining ports. Hence $L_0E_e$ constrains only block $u$ and
port $i$, while $L_0E_f$ constrains only block $w$ and port $j$.
These sets of tensor factors are disjoint. Their orthogonal projectors
$p,q$ therefore commute. This is the exact point where the simple-graph
hypothesis excludes the two outside endpoints from coinciding.

Put $M=LL_0^{-1}$. Its Gram operator on the central block is
\[
 G=M^*M=G_0^{-1/2}(\Phi_i\otimes\Phi_j\otimes\mathcal R)(C_v)G_0^{-1/2}
      =I+A+B+D.
\]
Its polar decomposition is $M=UG^{1/2}$ with $U$ unitary. Therefore the
angle between $LE_e,LE_f$ equals that between
$G^{1/2}\operatorname{ran}p,G^{1/2}\operatorname{ran}q$.
The single-defect term $A$ has identity on selected port $j$ after
whitening, and $B$ has identity on port $i$. Neither acts on the outside
endpoint blocks. Thus $[A,q]=[B,p]=0$, even though the retained central
ports may be mutually entangled.

The single CP sandwiches, tensoring with the actual CP map $\mathcal R$,
application to $C_v\ge0$, and congruence by $G_0^{-1/2}$ give
$-xI\le A\le xI$, $-yI\le B\le yI$. For Hermitian matrices these imply
the required operator-norm bounds. (CP-MIX) gives $-xyI\le D\le xyI$.
Tensoring the lower bounds gives $G\ge(1-x)(1-y)I$. Thus all hypotheses
of (SCHUR-FINAL) have now been produced from physical data.

\subsection{Producing \texttt{hcross}, including complements}
Let $P,Q$ be orthogonal projections. Diagonalize $PQP$ on
$\operatorname{ran}P$. An eigenvector $a$ of eigenvalue $c^2\in(0,1)$
produces the orthonormal partner
\[
 b=\frac{Qa-c^2a}{c\sqrt{1-c^2}}.
\]
Direct use of $Q^2=Q$ gives, on their span,
\[
 P=\begin{pmatrix}1&0\\0&0\end{pmatrix},\qquad
 Q=\begin{pmatrix}c^2&cs\\cs&s^2\end{pmatrix},\quad
 s=\sqrt{1-c^2}.
\]
Partners from orthogonal eigenvectors are orthogonal; the remaining
summands are commuting $0/1$ blocks. Eigenvalue one corresponds precisely
to the common range, which is removed in the Friedrichs cosine.
For $\xi=(a_0,b_0)$ on a nontrivial block, set $t=sa_0-cb_0$. Then
$(I-P)\xi=(0,b_0)$ and $(I-Q)\xi=(s,-c)t$, so
\[
 \operatorname{Re}\langle(I-P)\xi,(I-Q)\xi\rangle
 \ge-c|b_0||t|.
\]
The commuting blocks contribute nonnegative cross terms. Summing and
applying Cauchy--Schwarz to the block norms proves
\[
 \operatorname{Re}\langle(I-P)\xi,(I-Q)\xi\rangle
 \ge-c_F(P,Q)\|(I-P)\xi\|\|(I-Q)\xi\|.
\tag{HCROSS}
\]
Thus the preceding physical construction produces the swarm's
\texttt{hcross} with $C_{ef}=\mathcal B(\delta_{v,e},\delta_{v,f})$;
disjoint edge constraints commute and use $C_{ef}=0$.

Finally, every block-image projection commutes with every regional
projector: an incident regional image is contained in that block image,
and a nonincident operator acts on disjoint factors. On a sector outside
at least one block image, an incident complement term equals the
identity, so the energy is at least one. Every block has an incident
internal edge under the specified region hypotheses. Therefore the
matrix-gap theorem applies on the block-image sector and the same bound
extends to the full physical Hilbert space.

\paragraph{Formal implementation boundary.}
The reusable infrastructure needed is finite-dimensional spectral theory,
Kraus/Choi equivalence, tensor-product Hilbert spaces, orthogonal projections,
and positive square roots. The proof above derives the application-specific
injectivity, Schur metric, physical-subspace identification and angle bound.
None should be replaced by an assumed \texttt{hcross}.
